\documentclass[11pt]{article}
\usepackage[margin=1in]{geometry}
\usepackage{amsmath,amssymb,amsthm,mathtools}
\usepackage{microtype}
\usepackage[hidelinks]{hyperref}

\newtheorem{theorem}{Theorem}
\newcommand{\Sym}{\mathbb S^n}
\newcommand{\Frob}{\mathrm F}
\newcommand{\op}{\mathrm{op}}
\newcommand{\Id}{\mathrm{Id}}
\newcommand{\Dset}{\mathcal D}
\newcommand{\Qset}{\mathcal Q}
\newcommand{\Zset}{\mathcal Z}

\title{The exceptional Douglas--Rachford locus is porous}
\author{Autonomous mathematics run \texttt{fa\_banach\_001}, lane 17}
\date{13 August 2026}

\begin{document}
\maketitle

\begin{center}
\fbox{\parbox{0.91\textwidth}{
\textbf{Status: literature-implied answer (full scope), likely valid,
pending human review.}  The supporting polynomial-avoidance theorem does not
mention the source question; the identification below gives a complete
affirmative answer.
}}
\end{center}

\section{Source question and supporting theorem}

Let $X=\mathbb R^n$, $n\ge2$.  Bauschke, Schaad, and Wang
\cite{BSW} consider the space $\mathcal L\cap\mathcal S$ of symmetric
maximally monotone linear relations.  For $A$ in this space, write
\[
 J_A=(\Id+A)^{-1},\qquad R_A=2J_A-\Id.
\]
For a pair $(A,B)$, its Douglas--Rachford map is
\[
 T_{A,B}=\tfrac12(\Id+R_BR_A).
\]
They equip pairs with the metric
\begin{equation}\label{eq:source-metric}
 d((A,B),(A',B'))
 =\|J_A-J_{A'}\|_{\op}+\|J_B-J_{B'}\|_{\op}
\end{equation}
and prove that
\[
 \Dset=\{(A,B):T_{A,B}\text{ is a proximal mapping}\}
\]
is closed and nowhere dense.  Remark 3.3(ii), printed/PDF page 7 of
arXiv:1602.05626v1, asks:
\begin{quote}
``The set $D$ in Theorem 3.1 is closed and nowhere dense. Is it also a
porous set?''
\end{quote}

The later result used here is Corollary 1.7, printed/PDF page 3 of
Glazyrin--Karasev--Polyanskii, arXiv:2112.05382v5 \cite{GKP}:
if $P$ is a nonzero real polynomial of degree $m$, then the Euclidean unit
ball contains a point whose distance from $P^{-1}(0)$ is at least $1/m$.
By translation and rescaling, the same assertion holds in every Euclidean
ball, with the distance scaled by its radius.

\section{Identification and affirmative answer}

Let
\[
 K=\{R\in\Sym:\|R\|_{\op}\le1\},\qquad \Qset=K\times K.
\]
Reflected-resolvent coordinates give a bijection
\begin{equation}\label{eq:phi}
 \Phi:(\mathcal L\cap\mathcal S)^2\longrightarrow\Qset,
 \qquad \Phi(A,B)=(R_A,R_B).
\end{equation}
This includes singular boundary points: conversely, for $R\in K$ the
symmetric firmly nonexpansive matrix $(\Id+R)/2$ is the resolvent of a
symmetric maximally monotone linear relation.  If $\Qset$ is given the metric
\begin{equation}\label{eq:rho}
 \rho((R,S),(R',S'))
 =\tfrac12\bigl(\|R-R'\|_{\op}+\|S-S'\|_{\op}\bigr),
\end{equation}
then \eqref{eq:source-metric} shows that $\Phi$ is an isometry.

The source paper proves that a linear firmly nonexpansive map is proximal
exactly when it is symmetric.  Hence
\[
 T_{A,B}=T_{A,B}^{\mathsf T}
 \quad\Longleftrightarrow\quad
 R_BR_A=R_AR_B.
\]
It follows that
\begin{equation}\label{eq:commuting}
 \Phi(\Dset)=\Zset\cap\Qset,
 \qquad
 \Zset=F^{-1}(0),\quad
 F(R,S)=\|RS-SR\|_{\Frob}^2.
\end{equation}
Here $F$ is a nonzero real polynomial of degree four.  Nonzeroness is already
witnessed in the leading $2\times2$ block by
\[
 R=\begin{pmatrix}1&0\\0&-1\end{pmatrix},\qquad
 S=\begin{pmatrix}0&1\\1&0\end{pmatrix}.
\]

\begin{theorem}\label{thm:main}
For every $n\ge2$, the set $\Dset$ in Theorem 3.1 of \cite{BSW} is porous
in the metric \eqref{eq:source-metric}.  More quantitatively, at every point
of $\Dset$ and at every sufficiently small radius $r$, the ambient
$d$-ball of radius $r$ contains a $d$-ball of radius $r/(64n)$ disjoint from
$\Dset$.
\end{theorem}

\begin{proof}
Work first in the Euclidean space
$E=\Sym\times\Sym$ with product Frobenius norm
\[
 \|(R,S)\|_E=(\|R\|_{\Frob}^2+\|S\|_{\Frob}^2)^{1/2}.
\]
Fix $x=(R,S)\in\Zset\cap\Qset$, and let $r>0$ be small.  Put
\[
 \eta=\frac{r}{4\sqrt{2n}},\qquad x_0=(1-\eta)x.
\]
Because $R$ and $S$ commute, $x_0\in\Zset$.  Moreover,
\[
 \|x-x_0\|_E\le\eta\sqrt{2n}=r/4.
\]
The Euclidean ball $B_E(x_0,\eta)$ lies in $\Qset$: each component of
$x_0$ has operator norm at most $1-\eta$, while Frobenius norm dominates
operator norm.

Apply Corollary 1.7 of \cite{GKP} to
\[
 u\longmapsto F\bigl(x_0+(\eta/2)u\bigr).
\]
This is again a nonzero degree-four polynomial.  There is a point
$y\in B_E(x_0,\eta/2)$ satisfying
\[
 \operatorname{dist}_E(y,\Zset)\ge\eta/8.
\]
Consequently,
\[
 B_E(y,\eta/8)
 \subset B_E(x_0,\eta)\setminus\Zset
 \subset\Qset\setminus(\Zset\cap\Qset).
\]
It also lies in $B_E(x,r)$, because
\[
 r/4+\eta/2+\eta/8<r.
\]
Thus $\Zset\cap\Qset$ is uniformly porous in the Euclidean metric, with
hole ratio
\begin{equation}\label{eq:e-constant}
 c_E=\frac{1}{32\sqrt{2n}}.
\end{equation}

For $h=(H_1,H_2)\in E$, the norm comparisons are
\begin{equation}\label{eq:norms}
 \frac{\|h\|_E}{2\sqrt n}
 \le \rho(h)
 \le \frac{\|h\|_E}{\sqrt2}.
\end{equation}
Given a $\rho$-ball of radius $r$, apply the Euclidean construction inside
the $E$-ball of radius $\sqrt2r$, which is contained in that $\rho$-ball.
The Euclidean hole contains the $\rho$-ball with the same center and radius
\[
 \frac{c_E\sqrt2r}{2\sqrt n}=\frac{r}{64n}.
\]
This proves the asserted porosity of $\Zset\cap\Qset$ in $(\Qset,\rho)$.
The isometry \eqref{eq:phi} transfers it to $\Dset$.
\end{proof}

\section{Provenance, scope, and review}

This is a \emph{literature-implied} answer.  The authors of \cite{GKP} do
not cite \cite{BSW} or state an application to Douglas--Rachford maps.  The
implication uses the agent-identified commuting-locus reformulation
\eqref{eq:commuting}, plus the radial step needed because the parameter space
is the closed contraction ball rather than all of $E$.

The conclusion answers Remark 3.3(ii) in full for every finite dimension
$n\ge2$.  It does not address Remark 3.3(i), concerning extensions from
symmetric linear relations to general subdifferential operators.

The bounded search checked the run's solution, attempt, proof-gap, and ledger
indexes; the parsed arXiv corpus; and exact-title/exact-phrase and close-term
arXiv web searches using arXiv id 1602.05626, ``Douglas--Rachford'',
``proximal'', ``porous'', and ``commuting pairs''.  No later paper explicitly
answering Remark 3.3(ii) was found.  The supporting theorem was located by a
separate search for polynomial zero sets and porosity.

Human review should prioritize the boundary-point surjectivity of
\eqref{eq:phi} and the norm transfer \eqref{eq:norms}.  No computation is used
in the proof.

\begin{thebibliography}{9}
\bibitem{BSW}
H.~H. Bauschke, J.~Schaad, and X.~Wang,
\emph{On Douglas--Rachford operators that fail to be proximal mappings},
arXiv:1602.05626v1 (2016).

\bibitem{GKP}
A.~Glazyrin, R.~Karasev, and A.~Polyanskii,
\emph{Covering by planks and avoiding zeros of polynomials},
arXiv:2112.05382v5 (2022), Corollary 1.7.
\end{thebibliography}

\end{document}
